\scanpage{18}
inégalités bizarres dans les \unsure{énoncés} avec un terme (\uline{\unsure{signe}}).

\uline{Corollaire 1} \deleted{Soit $W$ une fonction de Weyl sur $(L^\infty)^+(\unsure{]0,\infty[})$ \unsure{et}} Soient $f$ \inserted{et} $g$ deux fonctions positives \deleted{décroissantes sur $(0,\infty)$} \inserted{sur l'espace mesuré $M$, et $N$,} telles que \deleted{$\Delta_t(f)<\unsure{0}$}
\[
(4)\qquad\Delta_t(f)\leq\Delta_t(g)<+\infty\qquad0\leq t<\omega
\]
Alors on a, pour toute fonction de Weyl \deleted{croissante $W$ sur} croissante sur $(L^\infty)^+(0,\infty)$ :
\[
(5)\qquad W(f)\leq W(g).
\]

On se ramène au cas où \deleted{$f$ et $g$ sont \unsure{à valeurs finies}, et} \inserted{$f$ et $g$ sont décroissantes sur $M=]0,\omega[$, avec $0<\omega<+\infty$, puis au cas où} $f$ et $g$ sont de la forme $\exp h$, et on applique la proposition \unsure{4}. \deleted{\unsure{La réciproque est encore vraie}} \inserted{On} peut faire des remarques analogues à celles qui suivent la prop. 4.

\deleted{\uline{\unsure{Corollaire}}}
\mspage{19}
